% problem_id: S001-more-morphisms-theorem-stein-factorization-Noetherian
% source_id: S001 (Stacks Project)
% source_file: more-morphisms.tex
% statement_locator: more-morphisms.tex lines 15204--15224
% proof_locator: more-morphisms.tex lines 15226--15287
% env: theorem
% id_note: label 跨文件重名 ⇒ id 用文件限定形式
% status: complete (statement + author proof verbatim + reason + locators)
%
\begin{theorem}[Stein factorization; Noetherian case]
\label{theorem-stein-factorization-Noetherian}
Let $S$ be a locally Noetherian scheme.
Let $f : X \to S$ be a proper morphism.
There exists a factorization
$$
\xymatrix{
X \ar[rr]_{f'} \ar[rd]_f & & S' \ar[dl]^\pi \\
& S &
}
$$
with the following properties:
\begin{enumerate}
\item the morphism $f'$ is proper with geometrically connected fibres,
\item the morphism $\pi : S' \to S$ is finite,
\item we have $f'_*\mathcal{O}_X = \mathcal{O}_{S'}$,
\item we have $S' = \underline{\Spec}_S(f_*\mathcal{O}_X)$, and
\item $S'$ is the normalization of $S$ in $X$, see
Morphisms, Definition \ref{morphisms-definition-normalization-X-in-Y}.
\end{enumerate}
\end{theorem}

\begin{proof}
Let $f = \pi \circ f'$ be the factorization of
Lemma \ref{lemma-stein-universally-closed}. Note that besides the
conclusions of Lemma \ref{lemma-stein-universally-closed} we
also have that $f'$ is separated
(Schemes, Lemma \ref{schemes-lemma-compose-after-separated})
and finite type
(Morphisms, Lemma \ref{morphisms-lemma-permanence-finite-type}).
Hence $f'$ is proper. By
Cohomology of Schemes, Proposition
\ref{coherent-proposition-proper-pushforward-coherent}
we see that $f_*\mathcal{O}_X$ is a coherent $\mathcal{O}_S$-module.
Hence we see that $\pi$ is finite, i.e., (2) holds.

\medskip\noindent
This proves all but the most interesting assertion, namely that
all the fibres of $f'$ are geometrically connected.
It is clear from the discussion above that we may replace $S$ by $S'$,
and we may therefore assume that $S$ is Noetherian, affine,
$f : X \to S$ is proper, and $f_*\mathcal{O}_X = \mathcal{O}_S$.
Let $s \in S$ be a point of $S$. We have to show that $X_s$ is
geometrically connected. By Lemma
\ref{lemma-characterize-geometrically-connected-fibres}
we see that it suffices to show $X_u$ is connected
for every \'etale neighbourhood $(U, u) \to (S, s)$.
We may assume $U$ is affine. Thus $U$ is Noetherian
(Morphisms, Lemma \ref{morphisms-lemma-finite-type-noetherian}),
the base change $f_U : X_U \to U$ is proper
(Morphisms, Lemma \ref{morphisms-lemma-base-change-proper}),
and that also $(f_U)_*\mathcal{O}_{X_U} = \mathcal{O}_U$
(Cohomology of Schemes, Lemma \ref{coherent-lemma-flat-base-change-cohomology}).
Hence after replacing
$(f : X \to S, s)$ by the base change $(f_U : X_U \to U, u)$
it suffices to prove that the fibre $X_s$ is connected when
$f_*\mathcal{O}_X = \mathcal{O}_S$. We can deduce this
from Derived Categories of Schemes, Lemma
\ref{perfect-lemma-proper-idempotent-on-fibre}
(by looking at idempotents in the structure sheaf of $X_s$)
but we will also give a direct argument below.

\medskip\noindent
Namely, we apply the theorem on formal functions,
more precisely Cohomology of Schemes, Lemma
\ref{coherent-lemma-formal-functions-stalk}.
It tells us that
$$
\mathcal{O}^\wedge_{S, s} = (f_*\mathcal{O}_X)_s^\wedge =
\lim_n H^0(X_n, \mathcal{O}_{X_n})
$$
where $X_n$ is the $n$th infinitesimal neighbourhood of $X_s$.
Since the underlying topological space of $X_n$ is equal to that
of $X_s$ we see that if $X_s = T_1 \amalg T_2$ is a disjoint union
of nonempty open and closed subschemes, then similarly
$X_n = T_{1, n} \amalg T_{2, n}$ for all $n$. And this in turn means
$H^0(X_n, \mathcal{O}_{X_n})$ contains a nontrivial idempotent $e_{1, n}$,
namely the function which is identically $1$ on $T_{1, n}$ and
identically $0$ on $T_{2, n}$. It is clear that $e_{1, n + 1}$
restricts to $e_{1, n}$ on $X_n$. Hence $e_1 = \lim e_{1, n}$
is a nontrivial idempotent of the limit. This contradicts the fact
that $\mathcal{O}^\wedge_{S, s}$ is a local ring. Thus the
assumption was wrong, i.e., $X_s$ is connected, and we win.
\end{proof}
